\subsection{Inference accounting and reuse thresholds}
\label{app:inference-amortization}
Basis-query timing follows the continuous evaluation boundary in \docref{Section}{sec:inference-cost}; numerical construction additionally charges geometry-specific mesh preparation, point mapping, factorization, and the spectral solve. \docref{Table}{tab:runtime-offline} reports all-node evaluation costs, while \docref{Appendix}{app:unseen-geometry} gives fresh numerical constructions on unseen domains. The reuse counts below concern trajectories that share spatial bases and operators.

\begin{table}[!htb]\centering\small
\caption{Reuse thresholds computed from original wall-time receipts, without rounding inputs. A reuse is one $T=1.16$, 1856-step trajectory on the same geometry, with the same time step and existing operators. The baseline is the numerical NS integration timer. Training-only counts set $T_{\rm off}=0$; adding recorded assembly or derived-field preparation charges those stages once. Derived-field costs include the recovered $K=512$ preparation timers. Per-query projection/output and separately unrecorded stages lie outside these integration-only thresholds.}
\label{tab:inference-amortization}
\begin{tabular}{rrrr}\toprule
$K$&Training only&+ $Q$ assembly&+ derived fields\\\midrule
512&123&154&179\\
1024&359&774&907\\
\bottomrule\end{tabular}\end{table}

For a full workload, replace the denominator of \doceqref{eq:inference-break-even} by the difference between complete per-query costs, including initial projection, selected field reconstruction, transfers, and storage. Output costs are assigned according to the requested number of reconstructed fields: the recorded NS workload saves 30 fields, and the ROM workload reconstructs three selected fields.

For basis generation over a geometry family, write $T_{\rm num}^{(g)}$ for a fresh numerical basis construction and $T_{\rm query}^{(g)}$ for the frozen neural map with the same postprocessing. For $G$ geometry queries indexed by $g$, training amortizes when
\begin{equation}
\sum_{g=1}^{G}\left(T_{\rm num}^{(g)}-T_{\rm query}^{(g)}\right)\geq T_{\rm train}.
\end{equation}
If the per-geometry costs are constant, write them as $T_{\rm num}$ and $T_{\rm query}$; then $G_\star=\lceil T_{\rm train}/(T_{\rm num}-T_{\rm query})\rceil$. This geometry-level threshold requires matched basis-generation timers. Each new geometry also incurs its own operator assembly when autonomous dynamics are requested. Arbitrarily dense neural spatial queries preserve the same learned continuous function; their cost scales with requested points and modes, and their accuracy is that of the learned representation.



\paragraph{Tensor-product and nonuniform parameter sweeps.}
A tensor-product sweep with $m$ values for each of $d$ independent query parameters contains $m^d$ cases. The per-axis break-even resolution $r_\star(d)$ and integer sampling count $m_\star(d)$ are
\begin{equation}
r_\star(d)=n_\star^{1/d},\qquad
m_\star(d)=\left\lceil r_\star(d)\right\rceil.
\label{eq:sweep-resolution}
\end{equation}
Here $d$ counts varied physical or query parameters and $n_\star$ is the integration-only reuse count in \doceqref{eq:inference-break-even}.


Unequal axis counts $m_1,\ldots,m_d$ meet the threshold when $\prod_{j=1}^d m_j\geq n_\star$; their geometric mean $(\prod_jm_j)^{1/d}$ is the mean per-axis resolution. Sparse, constrained, or correlated sweeps use the actual query count. With per-query savings $\Delta T_i$ and reusable cost $T_{\rm one\mbox{-}time}$, break-even requires $\sum_i\Delta T_i\geq T_{\rm one\mbox{-}time}$, including per-case preparation. Fixed-geometry initial-condition sweeps reuse the frozen basis and operators; geometry sweeps additionally charge geometry generation and required operator assembly.
